\documentclass[11pt]{article}
\usepackage[margin=1in]{geometry}
\usepackage{amsmath,amssymb,amsthm,hyperref,enumitem}
\hypersetup{colorlinks=true,linkcolor=blue,urlcolor=blue}
\newtheorem{theorem}{Theorem}
\newcommand{\rev}{\mathcal R}
\title{The algebraic Robbins Hankel conjecture}
\author{Principia Math In Progress}
\date{22 September 2026}
\begin{document}
\maketitle
\noindent\textbf{MathDB 350292. A proof of Barry's Conjecture 32.}

\section*{Task 1: the full determinant identity}
For a formal power series \(D\) with \(D(0)=1\), define its
revert transform by
\[
 \rev(D)(x)=\frac{(tD(t))^{\langle-1\rangle}(x)}x,
 \qquad H_n(D)=\det([x^{j+k}]D)_{0\le j,k\le n}.
\]
The inverse is compositional. Set \(H_{-1}(D)=1\), and let
\[
 A_n=\prod_{k=0}^{n-1}\frac{(3k+1)!}{(n+k)!},\qquad A_0=1
\]
be the Robbins numbers. All series below are formal over
\(\mathbb Q\), extended to \(\mathbb Q(i)\) when needed.

Use the constant-term-one square root to define
\[
 F(t)=\frac{\sqrt{1+2t+5t^2}-1-t}{2t^2},\qquad
 A(t)=\frac1{1-t^2/(1+t^2F(t))}.                       \tag{1}
\]
Then \(F(0)=1\) and
\[
 t^2F(t)^2+(1+t)F(t)-1=0.                             \tag{2}
\]
Consequently \(F=1/(1+t+t^2F)\), so \(A(ix)\) is precisely
the normalized continued fraction in [1, Conjecture~32]:
\[
 A(ix)=\cfrac1{1+\cfrac{x^2}{1-\cfrac{x^2}{1+ix-
              \cfrac{x^2}{1+ix-\cdots}}}}.
\]

\begin{theorem}
For every integer \(n\ge0\),
\[
 \boxed{\displaystyle H_n\bigl(\rev(A(i\,\cdot))\bigr)=A_n.} \tag{3}
\]
Equivalently, if \(U=\rev(A)\), then
\(H_n(U)=(-1)^{n(n+1)/2}A_n\).
\end{theorem}

\section*{Reduction to a rational reversion}
Let
\[
 G(t)=\frac{1+t}{1+t-2t^2-t^3},\qquad V=\rev(G).
\]
We prove the exact identity
\[
                         U(x)=\frac1{1+x^2V(x)}.       \tag{4}
\]
Indeed, put \(u=xU(x)\), so that \(uA(u)=x\). From (1),
\[
 U=1-\frac{u^2}{1+u^2F(u)}.
\]
Since \(U=1-x^2+O(x^3)\), the series
\(V_0=(U^{-1}-1)/x^2\) has constant term one. The last equation
gives \(V_0=U/(1+u^2F(u))\), hence
\[
 F(u)=\frac{U/V_0-1}{u^2},\qquad
 u=\frac{x}{1+x^2V_0}.
\]
Substitution in (2), followed by clearing denominators, gives
\[
 x^3V_0^3+2x^2V_0^2+xV_0^2-xV_0+V_0-1=0,
\]
or
\[
 (1-V_0)(1+xV_0)-x^2V_0^2(2+xV_0)=0.                \tag{5}
\]
Temporary division by \(x^2\) is legitimate in the Laurent
series field; all resulting expressions are the stated power
series. Equation (5) is equivalent to
\(xV_0G(xV_0)=x\). The unique compositional inverse with linear
coefficient one therefore gives \(V_0=V\), proving (4).

\section*{Deleting one Hankel row and column}
For any series \(D\), its Hankel matrix has generating kernel
\[
 K_D(x,y)=\frac{xD(x)-yD(y)}{x-y}
         =\sum_{j,k\ge0}[t^{j+k}]D(t)\,x^jy^k.
\]
Identity (4) implies
\[
 \frac{K_U(x,y)}{U(x)U(y)}=1-xyK_V(x,y).             \tag{6}
\]
Multiplication of a matrix kernel by \(1/U(x)\) on the left
and \(1/U(y)\) on the right is congruence by the lower triangular
Toeplitz matrix of \(1/U\). Its diagonal entries are one.
For each finite leading principal block, (6) therefore yields
\[
                  H_n(U)=(-1)^n H_{n-1}(V).          \tag{7}
\]
This also holds at \(n=0\) with the empty determinant convention.

\section*{The established rational input and the signs}
Barry's later paper [2, Proposition~9] proves the known formula
\[
 J(t)=\frac{1-t}{1-2t-t^2+t^3}
 \quad\Longrightarrow\quad
 H_m(\rev(J))=(-1)^{m(m+1)/2}A_{m+1}.                \tag{8}
\]
We explicitly credit this rational-series evaluation. Its
application to \(G\) can also be read in [2, Proposition~14].
For completeness the required invariance is proved here.
Put \(B(t)=G(-t)\) and \(W=\rev(B)\). Direct algebra gives
\(J=B/(1-tB)\). Inverting the functions \(tJ(t)\) and
\(tB(t)\) consequently gives
\[
 \rev(J)(x)=\frac1{1+x}W\!\left(\frac{x}{1+x}\right). \tag{9}
\]
Its coefficient of \(x^k\) is
\(\sum_{j=0}^k\binom kj(-1)^{k-j}[x^j]W\).
The resulting Hankel matrix is the congruence of the Hankel
matrix of \(W\) by the unipotent Pascal matrix
\(P_{kj}=\binom kj(-1)^{k-j}\).
Thus these Hankel determinants agree. Also \(V(x)=W(-x)\);
scaling coefficient \(k\) by \((-1)^k\) multiplies a size
\(m+1\) Hankel determinant by \((-1)^{m(m+1)}=1\).
We conclude from (8) that
\[
 H_m(V)=(-1)^{m(m+1)/2}A_{m+1}.
\]
Combining this with (7) proves
\(H_n(U)=(-1)^{n(n+1)/2}A_n\).
Finally, compositional inversion commutes with coefficient
scaling: \(\rev(A(c\,\cdot))(x)=U(cx)\) for \(c\ne0\).
At \(c=i\), the determinant multiplier is
\(i^{n(n+1)}=(-1)^{n(n+1)/2}\). The two signs cancel,
giving (3) at every \(n\ge0\).

The first real coefficients are
\(1,0,-1,0,3,-1,-12,11,51,-89,-204,628,646\), as in [1].
The proof uses the normalized definition (1) and the exact
Conjecture~32 continued fraction. The stray numerator \(x\)
in the preceding real continued-fraction display in [1] is
not used as grounds for any counterexample or research credit.

\section*{Tasks 2--5: proposed mathematical follow-ups}
\begin{enumerate}[label=\textbf{Task \arabic*:},start=2,leftmargin=*]
\item Classify the parameters in quadratic tails
\(at^2F^2+(1+bt)F-1=0\) and transforms
\(A=1/(1-ct^2/(1+dt^2F))\) for which the reduced series
\((\rev(A)^{-1}-1)/x^2\) is the reversion of a rational function.
\item Seek a weighted extension of (3) whose determinants
enumerate alternating sign matrices by their number of
negative entries, and determine the corresponding tail equation.
\item Determine explicit three-term recurrence coefficients
for the formal orthogonal polynomials associated with \(U\),
including the coefficients not determined by the Hankel products.
\item Construct a direct cancellation or nonintersecting-path
proof of (3) from the continued fraction, explaining why the
surviving objects are counted by \(A_n\).
\end{enumerate}

\begin{thebibliography}{9}
\bibitem{pre}
P.~Barry, \emph{Centered polygon numbers, heptagons and nonagons,
and the Robbins numbers},
\href{https://arxiv.org/abs/2104.01644}{arXiv:2104.01644v2}
(7 April 2021), Section~10, pp.~40--41, Conjecture~32.
\bibitem{final}
P.~Barry, \emph{From Fibonacci to Robbins: Series Reversion
and Hankel Transforms}, Journal of Integer Sequences
\textbf{24} (2021), Article 21.10.2,
\href{https://cs.uwaterloo.ca/journals/JIS/VOL24/Barry2/barry461.pdf}
{Propositions~9 and~14, pp.~8--11}.
This is the shortened and updated publication of [1];
the algebraic-series Conjecture~32 is not carried over.
\end{thebibliography}
\end{document}
